\section{Introduction}

LLM applications are increasingly written as programs rather than as single
prompts. A request is decomposed into a sequence of specialized agents, each of
which submits its own prompt to a shared serving engine, and each of those
prompts repeats a large body of content --- system instructions, tool
definitions, few-shot examples, accumulated context --- that a previous call
already paid to encode. Modern engines exploit this with prefix
caching~\cite{kwon2023pagedattention,zheng2023sglang}, but a cache only helps if
the right entry is resident at the right moment, and agentic traffic is
adversarial to recency-based replacement: the natural reuse distance of an
agent's context spans the entire rest of the workflow.

A line of recent systems responds by trying to know the future. KVFlow abstracts
the execution schedule as an agent step graph and prefetches the tensors an agent
will need in the next step~\cite{pan2025kvflow}. PBKV predicts agent invocations
several steps ahead by fusing historical workflows with the context of the target
workflow~\cite{zheng2026pbkv}. CacheScout learns agent execution transitions
online, with no predefined workflow graph and no offline training, and issues a
lightweight synthetic warmup request to materialize the predicted agent's
reusable anchor~\cite{zhang2026cachescout}. Pythia goes furthest: it mines
workflow structure from logged execution traces and prefills not only shared
prefixes but argument content recorded in earlier
executions~\cite{yu2026pythia}. All four place the same bet: that what the
program will do next can be reconstructed from outside the program. Helium
moves the knowledge into a compilation step instead --- and obtains it by
having the developer write the workflow down as a query plan in a dedicated
DSL~\cite{helium2026}.

\paragraph{Observation.} The bet is only necessary because the serving layer is
downstream of the program, and the DSL is only necessary because the host
language is assumed to be opaque. In a language-integrated LLM abstraction --- we build
on byLLM, the runtime of the meaning-typed programming model implemented in the
Jac language~\cite{dantanarayana2025mtp} --- the prompt is not handed to the
runtime by the application; it is \emph{assembled} by the compiler and runtime
from typed program constructs. A significant portion of what will eventually be
sent to the engine is therefore fixed at compile time. What is not fixed is
small and precisely characterized: the values that will be bound to the
parameters of each call. This changes the question from \emph{which agent runs
next?} to \emph{which prefixes are already knowable, and when does each argument
become available?} --- and both are answerable by static analysis.

\paragraph{Approach.} We build \sysname{}, a compiler--runtime co-design. The
compiler recovers a may-happen call topology as a BFS-layered reachable set,
recovers \emph{binding provenance} linking each produced value to the future
parameters it will feed, and separates the invariant from the variant part of
every prompt. Because it owns prompt assembly, it also rewrites the layout:
reusable zones are ordered first, and values that arrive late are confined to a
trailing binding zone, so that the reusable prefix is a genuine prefix rather
than being fragmented by interleaved dynamic content. The runtime consumes this
to prefill KV during engine idle windows --- both the temporal idle opened by a
tool call or a retrieval round trip and the spatial idle available under chunked
prefill~\cite{agrawal2024sarathi} --- constructing the longest knowable prefix:
invariants, graph context, already-bound argument values, and the grammar
compilation needed for constrained decoding~\cite{dong2025xgrammar}. Because the
served prompt remains byLLM's own rendering, speculation is advisory: a wrong
guess costs a cache miss, never a semantic change.

\paragraph{Contributions.}
\begin{enumerate}\itemsep2pt
\item A compiler-based static analysis that recovers, for byLLM programs and
      with zero annotation and zero training, the call topology, the provenance
      of each parameter, and the invariant portion of each runtime prompt, and
      that automatically lays out every LLM prompt for maximal KV reuse ---
      content that byLLM's stock runtime renders in a fragmented order.
\item A runtime that proactively and speculatively prefills KV cache in engine
      idle windows using that compile-time structure, reducing cold-start TTFT.
\item An evaluation of the resulting compiler--runtime co-design on two
      structurally different agent topologies; a broader study across agent
      applications is ongoing.
\end{enumerate}

\paragraph{Positioning.} We are not competing with KV residency management. KVFlow
and PBKV govern which KV \emph{stays} on the GPU; \sysname{} governs which KV
\emph{comes into existence}, and when. The two compose: a residency policy has
more to work with when the entries it is ranking were created ahead of demand
rather than on the critical path.
